\documentclass[fontsize=15pt]{jlreq}

%プリアンブル
\usepackage{amsmath}
\usepackage{mathtools}
\pagestyle{empty}
\usepackage[top=15mm, bottom=20mm, left=20mm, right=20mm]{geometry}

\newcommand{\m}[1]{\raisebox{0.1em}{\textcircled{\scriptsize #1}}}



%本文
\begin{document}
\begin{center}
{\LARGE {振り返り}} \\
-物理\m{2}- %単元ごとに変えるか消す
\end{center}
\vspace{1em}
\begin{flushright}
名前：\underline{\hspace{5cm}}
\end{flushright}

% 問題部分
\begin{enumerate}
    \item 等速直線運動となる条件をかけ．\\
    (\hspace{15cm})
    \vspace{2em} 

    \item 次のうち，だんだん速くなる運動を選べ．
    \begin{itemize}
        \item A：斜面を滑り落ちる運動
        \item B：水平面上で摩擦力がはたらく運動
        \item C：自由落下
        \item D：運動と逆向きにはたらく運動
    \end{itemize}
    (\hspace{4cm})
    \vspace{2em} 

    \item \m{1}だんだん速くなる運動と\m{2}等速直線運動において，移動距離と時間にはどのような関係があるか．\\
    \m{1} \\[0.5em]
    \m{2}
    \vspace{2em} 

    \item 物体を真上(鉛直)方向に投げ上げたとき，その物体はどのような運動をするか．\\
    \small ヒント：投げ上げたすぐと，少し時間が経ったあとでは違った運動をします．それぞれ2種類の運動を説明してください．\\
    回答：
    \vspace{2em} 

\end{enumerate}

\end{document}
